\documentclass[12pt]{article}

\input{../../_assets/preambulo.tex}


\begin{document}

    % 1. Foto de fondo
    % 2. Título
    % 3. Encabezado Izquierdo
    % 4. Color de fondo
    % 5. Coord x del titulo
    % 6. Coord y del titulo
    % 7. Fecha

    
    \input{../../_assets/portada}
    \portadaExamen{ffccA4.jpg}{Ecuaciones\\Diferenciales II\\Examen X}{Ecuaciones Diferenciales II. Examen X}{MidnightBlue}{-8}{28}{2026}{}

    \begin{description}
        \item[Asignatura] Ecuaciones Diferenciales II.
        \item[Curso Académico] 2024/2025.
        \item[Grado] Doble Grado en Ingeniería Informática y Matemáticas.
        \item[Grupo] Único.
        \item[Profesor] Rafael Ortega Ríos.
        \item[Descripción] Primera prueba de clase.
        \item[Fecha] 31 de Marzo de 2025.
        % \item[Duración] ---.
    
    \end{description}
    \newpage


    % ------------------------------------
    
    \begin{ejercicio}
        Dado un campo continuo $X : \left[0,+\infty\right[ \times \bb{R}^2 \to \bb{R}^2$, $(t,x) \mapsto X(t,x)$, y dado $\veps > 0$, demuestra que las fórmulas que siguen definen una función:
        \Func{x_\veps}{ \left[0,+\infty\right[}{\bb{R}^2}{t}{
            \begin{cases}
                (0,1), & t \in \left[0, \veps\right],\\
                x_\veps(t-\veps) + \veps X(t-\veps, x_\veps(t-\veps)), & t \in \left]\veps, +\infty\right[.
            \end{cases}
        }
        ¿Se puede asegurar que esta función es continua?
    \end{ejercicio}

    \begin{ejercicio}
        Demuestra que el problema de valores iniciales
        \[
            \dot{x}_i = \sqrt{1-t^2}(1-f(x_1,x_2))x_i, \quad x_i(0) = \frac{1}{2}, \quad i = 1,2
        \]
        admite una solución definida en $\left]-1,1\right[$, donde:
        \Func{f}{\bb{R}^2}{\bb{R}}{(x_1,x_2)}{\min(1, x_1^2 + x_2^2).}
    \end{ejercicio}

    \begin{ejercicio}
        Para cada $n \geq 1$ se considera la función
        \[
            x_n(t) = 1 + t + \frac{t^2}{2!} + \dots + \frac{t^n}{n!}, \quad t \in \left[-2,2\right].
        \]
        Encuentra una sucesión de números $\{\veps_n\}_{n \geq 1}$ convergente a cero para la que se cumpla
        \[
            \left|x_n(t) - 1 - \int_0^t x_n(s)ds\right| \leq \veps_n \quad \forall t \in \left[-2,2\right].
        \]
    \end{ejercicio}

    \begin{ejercicio}
        Se supone que $x : \left[-1,1\right] \to \bb{R}$ es una función continua que cumple, para cada $t \in \left[-1,1\right]$,
        \[
            x(t) = 2 + t \int_0^t x(s)ds.
        \]
        Encuentra un problema de valores iniciales para una ecuación diferencial de segundo orden del que $x(t)$ es solución. Deduce entonces que la ecuación integral admite a lo sumo una solución.
    \end{ejercicio}

    \begin{ejercicio}
        Nos dan una sucesión de funciones $\{f_n\}_{n \geq 0}$, $f_n : \left[0,1\right] \to \bb{R}$, que cumplen
        \begin{enumerate}[label=\roman*)]
            \item $|f_n(0)| \leq 7$, para cada $n = 0, 1, 2, \dots$
            \item $\{f_n\}_{n \geq 0}$ es equicontinua.
        \end{enumerate}
        ¿Se puede asegurar que existe una sucesión parcial de $\{f_n\}_{n \geq 0}$ que converge uniformemente en $\left[0,1\right]$?
    \end{ejercicio}

    \newpage
    \setcounter{ejercicio}{0} % Reiniciar contador de ejercicios
    \noindent
    % \textbf{Solución.}

\end{document}